% 精读章: Deschênes & Greenstone, The Economic Impacts of Climate Change:
%         Evidence from Agricultural Output and Random Fluctuations in Weather
%         (AER 2007，97(1) 三月号正刊)
% 原文出处: doi:10.1257/aer.97.1.354（所据版本：期刊已刊版 AER 97(1): 354--385；
%          本仓 papers/01_methods/2007-AER-Deschenes-Greenstone-*.pdf 为团队正典，
%          本章页码引注均指该版）
% 体例注记: 2026-10-21 起，正文中的 terms/expressions 一律用英文原词，
%          中文仅保留句法成分；定义处随文给中文对照并登记英文 \index。
% 建档/更新日期由章首 \chdate 渲染（grep: chdate）
% 主笔: bjr；初稿: 2026-11-04（Claude 协助起草，bjr 审订）
% 核对: 2026-11-19 对照本地 PDF 全章重核——原句引用逐条改回原文字面；
%       改正 Table 2C、Table 3、SHF 复刻、Table 5/6/8、样本县数、置信区间、
%       分州显著性等处的数字与归属；式 (3)(4) 记号改从原文。
% partinfo 定位: 把 weather 波动面板（panel approach）引入农业影响评估，
%               并自认面板估计不含全部 substitution、相对 long-run effect
%               向下偏（adaptation 批评的源头），见
%               \chref{ch:info-climate-change-adaptation} 的方法线与速查表。

\chapter{The Economic Impacts of Climate Change（Deschênes \& Greenstone，AER 2007）}
\label{ch:deschenes-greenstone-2007}
% 章题过长，页眉用短题（preamble 的 \chaptermark 经 \markboth 排页眉）
\chaptermark{Economic Impacts of Climate Change（Deschênes \& Greenstone，2007）}
\chdate{2026-11-04}{2026-11-19}

\section{这篇论文在说什么}

一句话概括：论文用 counties 内部逐年 \index{weather}weather 波动（conditional on county 与 state$\times$year fixed effects）估计 climate change 对美国农业部门利润的影响，以此回应截面 hedonic（即 \index{Ricardian model}Ricardian）路线的 omitted variables 批评；其经验结论是影响小且不显著——``The preferred estimates indicate that climate change will increase annual profits by \$1.3 billion in 2002 dollars (2002\$) or 4 percent''\citep[第~354 页]{deschenes2007economic}，并判定截面路线``extremely sensitive to seemingly minor choices about control variables, sample, and weighting''\citep[第~354 页]{deschenes2007economic}，故不可靠。

它在 partinfo 里的坐标：panel approach 的引入站（\chref{ch:info-climate-change-adaptation} 方法线 panel 线第一站）。论文自己写明的局限——面板估计不含全部 adaptation、相对 long-run effect 向下偏（第~359 页）——就是后来 adaptation 检验路线的起点，由 \chref{ch:burke-hsiang-miguel-2015} 与 \chref{ch:merel-gammans-2021} 接手。

骨架表（所据版本：AER 97(1) 已刊版；页码为期刊页码）：

\begin{table}[htbp]
\centering\small
\begin{tabular}{@{}p{3.4cm}p{1.6cm}p{7.8cm}@{}}
\toprule
原文章节 & 页码 & 讲什么 \\
\midrule
Abstract & 354 & +\$1.3 billion（4\%）的基准结论；hedonic 不可靠 \\
引言（未编号） & 354--357 & production function、hedonic 两条旧路线的缺陷；论文设计与主要发现 \\
\S I Conceptual Framework & 357--360 & I.A 面板新策略与 hedonic 的差别；I.B 用年际波动推断 climate 效应的经济学（式 (1)(2)、Figure~2），自认 downward bias \\
\S II Data Sources and Summary Statistics & 360--365 & Census of Agriculture 县级面板；PRISM 与气象站数据、growing season degree-days 构造；climate 预测情景；weather variation 统计（Table~2） \\
\S III Econometric Strategy & 365--367 & III.A hedonic 设定（式 (3)）；III.B 面板设定（式 (4)）与 identification assumption \\
\S IV Results & 367--380 & IV.A hedonic 重估（Table~3--4）；IV.B 利润面板（Table~5--7）；IV.C 单产（Table~8） \\
\S V Interpretation & 380--381 & 福利换算（现值 \$26 billion）；四条 caveats，另加两条会使损害被高估的 caveats \\
\S VI Conclusions & 381 & 结论与后续研究方向 \\
\bottomrule
\end{tabular}
\end{table}

\section{预备知识}

读者需备两件旧工具。其一，论文之前的两条路线各有一难：\index{production function approach}production function approach（生产函数法）估计 weather 对具体作物单产的影响，剔除了土壤等农户控制不了的因素，但``farmer adaptations are completely constrained in the production function approach''\citep[第~355 页]{deschenes2007economic}，估计向下偏；\index{hedonic approach}hedonic approach（特征价格法，即 Ricardian 路线）用农地价值对 climate 变量回归，理论上地价``reflect the present discounted value of land rents into the infinite future''\citep[第~355 页]{deschenes2007economic}，``In principle, this approach accounts for the full range of farmer adaptations''\citep[第~355 页]{deschenes2007economic}，但 unmeasured characteristics（土壤质量、转作他用的期权价值等，作者引 Hoch 1958、1962 与 Mundlak 1961）与 climate 相关，``the sign and magnitude of the resulting omitted variables bias is unknown''\citep[第~355 页]{deschenes2007economic}。其二，\index{degree-days}degree-days（度日，又称生长度日；术语来自农学与气象学）：作物累积热量暴露的度量，日平均温度在基点与顶点之间时按超出基点的度数累积（构造口径见下文数据节）。

\section{数据与测度}

\textbf{产出与利润。}农业数据来自 Census of Agriculture（1978、1982、1987、1992、1997、2002 六次普查），观测单位是县；普查的填报口径是当年农产品销售额不低于 \$1{,}000 的农场与牧场（第~360 页）。生产支出在 1978 与 1982 年没有采集，所以利润分析以 1987--2002 四次普查为基础（第~360 页）。样本有三个口径，引用时要分清：Table~1 的描述统计用 1978--2002 的平衡面板，2{,}268 县（第~362 页及脚注~11；Alaska 与 Hawaii 剔除，土壤、农地面积、灌溉面积、人均收入、人口密度、纬度任一变量缺失的县剔除）；利润回归再要求利润变量在各普查年都不缺失，得 2{,}262 县的平衡样本（第~372 页脚注~23）；hedonic 重估（Table~4）的样本又少 144 个县（第~371 页）。被解释变量是县级每英亩农地的利润：农产品销售的市场价值减去生产支出，\textbf{不含联邦农业项目支付}（第~360 页）。作者明言这是不理想的口径：普查的支出项不含未抵押自有资本的利息，只部分计入了农场的资本成本（第~360 页），利润因而偏高；全县平均地租约 \$35（2002\$），利润约 \$42（2002\$），因此把利润口径的影响估计乘以二者之比 0.83 才换算成租金口径的福利度量（第~360 页）。

\textbf{climate 与 weather。}climate 数据用 PRISM 模型（4$\times$4 公里格点的降水与温度插值，作者称其被``almost all professional weather services... regarded as one of the most reliable interpolation procedures''\citep[第~361 页]{deschenes2007economic}）。growing season（生长季）取 4 月 1 日至 9 月 30 日，覆盖冬小麦以外的多数作物（第~361 页）。degree-days 的构造遵循农学传统（Ritchie 与 NeSmith 1991）：基点 46.4\,$^\circ$F（8\,$^\circ$C）、顶点 89.6\,$^\circ$F（32\,$^\circ$C），日平均温度在区间内贡献（均温$-$基点）个 degree-days，超过顶点的日子贡献 43.2 个 degree-days（第~361--362 页）。Ritchie 与 NeSmith 还讨论过 93.2\,$^\circ$F（34\,$^\circ$C）处可能存在一个有害温度阈值，超过它升温即造成损害；作者在 robustness check 里考察了这一点（第~361 页；即 Table~6 的 harmful degree-days 一行，第~376 页）。月度 PRISM 数据算不出逐日 degree-days，作者改用约 8{,}000 个地面气象站的逐日数据（NCDC Cooperative Summary of the Day）——这是对既有文献（用月度数据加分布假设，如 Thom 1966 与 SHF 2006）的测度改进（第~362 页）。

\textbf{climate 情景。}两组预测：基准情景（全季全区均匀升温 5\,$^\circ$F、降水 +8\%）；Hadley 2 模型的逐日州级预测（2020--2049 与 2070--2099 两窗口，第~362 页）。

\textbf{weather 波动与 climate change 的量级对比。}这是测度一节最该记住的表。Table~2C（第~364 页）报告去年份 fixed effects、去 state$\times$year fixed effects 之后，县年观测的偏离超过各幅度的占比（1987、1992、1997、2002 四年平均）。去 state$\times$year 效应后，24.5\%、9.3\%、2.2\% 的县年 degree-days 偏离超过 400、800、1{,}200；62.3\%、35.3\%、18.1\% 的县年 growing season 降水偏离超过 1.0、2.0、3.0 英寸（第~365 页）。把它与要外推的幅度对比：Hadley 2 长期情景的平均增幅约 1{,}200 degree-days（Table~2B；基准情景的增幅约 720 degree-days，Table~2A），这样大的 degree-days 偏离在数据里只来自 2.2\% 的观测，作者的措辞相应有所保留（幅度约 1{,}000 degree-days 的偏离占 5\%）；降水则宽裕得多，超过 18\% 的观测偏离大于 3.0 英寸（第~365 页）。结论是 identification 所用的 variation 与要外推的 variation 在量级上有重叠，降水的重叠充分，degree-days 的重叠落在分布的尾部。这是 panel approach 可用于外推的前提，也是它的边界。

\section{模型设定与识别策略}

本节是论文的理论贡献：先把截面 hedonic 写出来（式 (3)），再给出面板替代（式 (4)），由 identification 反推 assumption，最后用 \S I.B 的经济学说明这套估计量测到的是什么、偏向哪边。记号随原文：$\overline{W}_{ic}$ 是县 $c$ 的第 $i$ 个 climate 变量（长期均值，原文称 climate normals），$W_{ict}$ 是它在年 $t$ 的 weather 实现。

\subsection{截面 hedonic 及其问题}

hedonic 设定为（式 (3)，第~365 页）
\begin{equation}
y_{ct} = \mathbf{X}_{ct}'\boldsymbol{\beta} + \sum_i \theta_i f_i(\overline{W}_{ic}) + \epsilon_{ct},
\qquad \epsilon_{ct} = \alpha_c + u_{ct},
\label{eq:dg-hedonic}
\end{equation}
其中 $y_{ct}$ 是县 $c$ 在年 $t$ 的每英亩农地价值；climate 变量取 1 月、4 月、7 月、10 月的温度与降水共八个，这一选择沿 Mendelsohn、Nordhaus 与 Shaw（1994）；$f_i$ 取线性加二次项是文献惯例，作者写明各 climate 变量的恰当函数形式并不清楚（第~366 页）；灌溉县与非灌溉县参数分开，则是采纳 SHF（2005）的意见（第~366 页）。一致估计要求 $\E[f_i(\overline{W}_{ic})\,\epsilon_{ct}\mid \mathbf{X}_{ct}]=0$。误差项里既有不随时间变化的 $\alpha_c$，也有逐年变化的 $u_{ct}$，只要其中任何一部分未被度量且与 climate 相关，这个条件就不成立（第~366 页）。土壤质量、灌溉可得性、人口密度都是这样的因素；Table~3 的平衡性检验给出了直接的证据（见结果节）。

\subsection{面板替代与 identification}

新设定把 climate 变量换成逐年 weather（式 (4)，第~367 页）：
\begin{equation}
y_{ct} = \alpha_c + \gamma_t + \mathbf{X}_{ct}'\boldsymbol{\beta} + \sum_i \theta_i f_i(W_{ict}) + u_{ct},
\label{eq:dg-panel}
\end{equation}
被解释变量改为县级农业利润，$\alpha_c$ 为完整的 \index{county fixed effects}county fixed effects，偏好设定把年份效应 $\gamma_t$ 换成 \index{state-by-year fixed effects}state$\times$year fixed effects $\gamma_{st}$（第~367 页）。被解释变量为何从地价换成利润：``land values capitalize long-run characteristics of sites and, conditional on county fixed effects, annual realizations of weather should not affect land values. Weather does, however, affect farm revenues and expenditures, and their difference is equal to profits''\citep[第~367 页]{deschenes2007economic}。climate 变量为何必须退场：``Second, it is impossible to estimate the effect of the long-run climate averages in a model with county fixed effects, because there is no temporal variation in $\overline{W}_{ic}$. Consequently, we replace the climate variables with annual realizations of weather, $W_{ict}$''\citep[第~367 页]{deschenes2007economic}。

identification 依赖
\begin{equation*}
\E\bigl[f_i(W_{ict})\,u_{ct} \mid \mathbf{X}_{ct}, \alpha_c, \gamma_{st}\bigr] = 0 ,
\end{equation*}
即 weather 参数由 county-specific deviations in weather about the county averages（再调整掉州内各县共同的冲击）识别；作者称这一 variation``is presumed to be orthogonal to unobserved determinants of agricultural profits''\citep[第~357、367 页]{deschenes2007economic}。推断上，全文报告 Eicker--White 异方差稳健标准误，偏好的 fixed effects 设定另报告 Conley 空间相关标准误（经纬度各 7 度、约 500 英里截断；第~366 页及脚注~15）；利润方程以农地总面积的平方根加权（第~374 页）。

\begin{pitfall}[fixed effects 放大 measurement error]
\label{pit:dg-attenuation}
作者自己提醒：``A shortcoming of this approach is that all the fixed effects are likely to magnify the importance of misspecification due to measurement error, which generally attenuates the estimated parameters''\citep[第~367 页]{deschenes2007economic}。fixed effects 削去大部分 variation 之后，weather 变量的 measurement error（测量误差）在剩余 variation 中的占比上升，\index{attenuation bias}attenuation bias（衰减偏误，系数向零偏）随之加重。读 Table~5--6 的系数时必须带着这一条：面板系数偏小，未必是「影响确实小」。
\end{pitfall}

\subsection{理论贡献：这套估计量测到的是什么（式 (1)(2)）}

\S I.B 把「用年际波动推断 climate 效应」的经济学摊开。代表性农户生产给定作物、不能随 weather 换作物，利润为（式 (1)，第~358 页）
\begin{equation}
\pi = p(q(w))\,q(w) - c(q(w)),
\label{eq:dg-profit}
\end{equation}
其中 $q$ 是产出，$w$ 是 weather；价格与总成本是产量的函数，产量是 weather 的函数。对 $w$ 求导（式 (2)，第~358 页）：
\begin{equation}
\frac{\partial \pi}{\partial w} =
\underbrace{\frac{\partial p}{\partial q}\frac{\partial q}{\partial w}\,q}_{\text{价格效应}}
+ \underbrace{\Bigl(p - \frac{\partial c}{\partial q}\Bigr)\frac{\partial q}{\partial w}}_{\text{数量效应}} .
\label{eq:dg-decomp}
\end{equation}
第一项是 weather 冲击经由价格变化带来的利润变化（乘以初始产量 $q$）：减产推高价格，这在短期内缓和代表性农户的损失。第二项是论文想隔离的效应（holding constant crop choice，第~358 页）。关键在两种时间尺度下价格反应的对比：长期里其他农户可以扩产、新农户可以进入，``it is sensible to assume that $(\partial p/\partial q)_{\text{Long Run}} > (\partial p/\partial q)_{\text{Short Run}}$ and is perhaps even equal to zero''\citep[第~358 页]{deschenes2007economic}（$\partial p/\partial q<0$，「$>$」指长期的值更接近零）。所以第一项短期可为正，长期则小得多甚至为零（第~358 页）。state$\times$year fixed effects 的作用之一就是吸收价格的共同变动，使估计接近第二项。

于是得到论文最要紧的一段自我定位：``This paper's empirical strategy relies on year-to-year variation in weather, and thus it is unlikely that farmers are able to switch crops upon a year's weather realization. The import for the subsequent analysis is that our estimates of the impact of climate change may be downward biased, relative to the preferred long-run effect that allows for all economic substitutions. If the degree of climate change is `small', however, our estimates are equal to the preferred long-run effect''\citep[第~359 页]{deschenes2007economic}。Figure~2 的几何版本：作物 1 的利润函数在 $T_1$ 到 $T_2$ 间最优、作物 2 在 $T_2$ 到 $T_3$ 间最优，hedonic 均衡是两段的上包络；climate 从 $T_1$ 升到 $T_3$，能换作物的长期利润读在 $C$ 点，被锁定在作物 1 的短期利润读在 $C'$ 点，$C'-A$ 是 $C-A$ 的向下偏估计（第~359 页）。\textbf{这是 panel approach 与 long-run effect 之间差距的第一次明确陈述}（只给了方向，没有给大小），也是 \chref{ch:merel-gammans-2021} 全部问题的起点。

\begin{intuition}[一年的 weather 与一地的 climate：哪些调整来得及]
D-G 区分的是调整所需的时间。面对 climate 的变化，农户可以调整化肥用量、改变作物组合，甚至把农地改作他用（作者举的例子是住宅开发；第~355 页）；面对某一年的 weather 实现，换作物这类调整来不及做（第~359 页），但投入品的组合仍可调整，作者举的例子是化肥与灌溉用水（第~360 页），所以面板估计仍优于完全不容许 adaptation 的 production function 估计。年际 weather 波动识别的是「作物选择给定、投入可调」之下的响应。\chref{ch:merel-gammans-2021} 用灌渠与风扇的例子表达的是同一区分。
\end{intuition}

\section{结果与解释}

按论文自己的顺序：先检验 hedonic 的前提，再重估 hedonic，最后给出面板估计。

\subsection{平衡性检验：climate normals 与可观测变量不独立}

Table~3（第~368 页）把县按 climate normals 分四分位组，比较各组可观测变量的均值。以 7 月温度为例，每英亩农地价值在四组依次为 1{,}118、1{,}771、1{,}609、1{,}481 美元（$F$ 统计量 25.5），人口密度与人均收入的组间差异同样显著（$F$ 统计量 19.7 与 14.1），均拒绝各组均值相等。作者先说明为什么平衡性值得检验：若可观测变量在各组间平衡，就有两点理由相信 hedonic 估计可靠——``First, consistent inference will not depend on functional form assumptions on the relation between the observable confounders and farm values. Second, the unobservables may be more likely to be balanced''\citep[第~368 页]{deschenes2007economic}。检验结果是不平衡：扩展检验中，112 个解释变量均值相等的原假设有 111 个在 1\% 水平被拒（第~369 页）。作者特别指出人口密度（可视作城市化程度或农地转为住宅的可能性的度量）与农地价值相关，而它对农业单产并无直接影响，这削弱了 hedonic 路线的可信度（第~369 页）。

\subsection{hedonic 重估：不稳健}

复刻 SHF（2005）的设定与数据，基准情景（升温 5\,$^\circ$F、降水 +8\%）对农地价值的总影响为 $-\$543.7$ billion（cropland 权重）或 $+\$69.1$ billion（crop revenue 权重）（第~369 页）。作者再换两组协变量：只保留 climate 变量、去掉其余协变量，估计变为 $-\$98.5$ 与 $+\$437.6$ billion；在 SHF 设定上加 state fixed effects，得到 $-\$477.8$ 与 $+\$1{,}034$ billion——后一个数字``seems implausible, since it is nearly as large as the entire value of agricultural land and buildings in the United States, which was \$1,115 billion in 2002''\citep[第~369 页]{deschenes2007economic}。改用农地面积加权并重构 SHF 数据后，三种设定的估计依次为 $+\$225.1$、$-\$315.4$、$-\$0.6$ billion（第~369 页）。作者结论：``It seems reasonable to conclude that the application of the hedonic approach to the SHF data fails to produce robust estimates of the impact of climate change, even with a single year of data''\citep[第~369 页]{deschenes2007economic}。

D-G 自己的更完整 hedonic（Table~4，第~370--371 页）把同样的不稳定摆得更齐：42 个估计中，偏好的 B、C 两组设定下分年度的点估计从 $-\$202$ 到 $\$321$ billion；六种合并设定与权重组合下从 $-\$75.1$ 到 $\$110.8$ billion；D-G 偏好的 C 设定给出显著为正的 $\$95.2$ 与 $\$110.8$ billion（第~371 页）。作者的评价是：``These results are troubling, because there is no ex ante reason to believe that the estimates from a particular year are more reliable than those from other years''\citep[第~371 页]{deschenes2007economic}。\textbf{注意读法}：这一小节的主要发现是同一数据稍换口径结果就翻转，即不可靠本身（摘要：``to be unreliable because it produces estimates that are extremely sensitive to seemingly minor choices about control variables, sample, and weighting''\citep[第~354 页]{deschenes2007economic}）；作者也顺带指出，其中较可信的 hedonic 设定显示 climate change 对美国农业总体略有益处（第~372 页），但没有把它当作主要结论。

\subsection{论文的 panel approach：利润的影响小且不显著}

Table~5（第~375 页）报告三个情景（基准、Hadley 2 的 2020--2049、Hadley 2 的 2070--2099）下利润的预测变化。偏好设定（第 4 列：土壤控制变量加 state$\times$year fixed effects）在 Hadley 2 长期情景下的结果是``climate change will lead to an increase in agricultural sector profits of roughly \$1.3 billion. This estimate is statistically indistinguishable from zero with either the Eicker-White (parentheses) or Conley (square brackets) standard errors''\citep[第~374 页]{deschenes2007economic}；用较小的标准误，95\% 置信区间约为 $-\$0.5$ 到 $\$3.1$ billion（第~374 页），相对 1987--2002 年均利润 \$32.3 billion 即 $-1.5\%$ 到 $+9.6\%$（第~377 页）。分解上，影响几乎全部来自非灌溉县（第~374 页；Table~5 第 4 列：非灌溉县 1.10、灌溉县 0.18，单位 billion），按变量分则主要来自降水（degree-days 0.39、降水 0.89）。Table~6（第~376 页）共十二行 robustness check，包括线性、三次式等温度与降水的其他函数形式、harmful degree-days、剔除强影响点、灌溉县定义的分界改为 5\% 与 15\%（基准为 10\%）、growing season 的定义与不加权等，估计全部落在 $-\$0.98$ 到 $\$3.59$ billion 之间；对前三行的函数形式变体，作者写明``a zero effect cannot be rejected in any of the cases''\citep[第~376 页]{deschenes2007economic}。分州看（Table~7，第~377--378 页），最大的输家是 California（$-\$750$ million，约其农业利润的 15\%），最大的赢家是 South Dakota（$+\$720$ million）与 Georgia（$+\$540$ million）；分州估计总体不精确，只有八个州能在 5\% 水平拒绝零影响（第~377 页）。单产端（Table~8，第~379 页），偏好设定下玉米与大豆全国产量的预测变化为 $+0.01$ 与 $+0.02$（单位：十亿蒲式耳），都统计上不异于零；作者据此认为利润影响小不太可能只是短期价格上涨所致（第~380 页）。

贯穿本章的例子在这里收尾。数据节 Table~2C 的量级对比说明外推所需的偏离在历史数据中出现过（degree-days 落在尾部，降水宽裕，第~365 页）；引言对 Figure~1 的概括则给出响应面的形状：``The striking finding is that all of the response surfaces are flat over the ranges of the predicted changes in degree-days and inches. If anything, climate change appears to be slightly beneficial for profits and yields''\citep[第~355 页]{deschenes2007economic}。两者合起来就是基准结论的直观版本：在历史上观测得到的偏离范围内，利润与单产对 degree-days 和降水的响应是平的。

\begin{pitfall}[从「影响不显著」读出「没有影响」]
D-G 的结论是``economically and statistically small changes in profits''\citep[第~377 页]{deschenes2007economic}，配一个置信区间（$-\$0.5$ 到 $\$3.1$ billion，第~377 页），不是影响为零的点断言。同样不可忘的是易错点~\ref{pit:dg-attenuation}：fixed effects 放大 measurement error，估计偏向零。后续文献对论文的批评（见「局限与批评」节）恰恰不是「估出了零」，而是「这套 variation 只测到 short-run response」——两件事不要混。
\end{pitfall}

\section{机制与讨论}

福利换算：假定 +\$1.3 billion 的年度利润增量对未来所有年份成立，并取 5\% 的折现率，``This implies that climate change increases the present value of the stream of land rents by about \$26 billion. The 95-percent confidence interval is $-\$10$ billion to \$62 billion''\citep[第~380 页]{deschenes2007economic}。作者强调这个区间``is much tighter than the range of point estimates from the hedonic approach''\citep[第~380 页]{deschenes2007economic}——面板路线的收获是区间的可信性，不只是点估计。换算的两个前提（利润增量对所有未来年份成立、5\% 折现率）是原文正文写明的 assumption，不在表内；引用 \$26 billion 时应一并交代。

\section{局限与批评}

论文自己承认的（\S V，第~380--381 页）。先是四条一般性的 caveats：其一，极端事件（干旱、洪水）与 climate 变率的变化不在 growing season degree-days 与降水的刻画里，论文的分析对这类事件``uninformative''\citep[第~380 页]{deschenes2007economic}；其二，climate change 是全球性的，会改变长期生产成本与相对价格，而年际波动加 state$\times$year fixed effects 的估计不含这类价格变动（作者同时指出 hedonic 同样处理不了）；其三，政府补贴项目影响利润与地价（引 Kirwan 2005）；其四，利润口径低估资本成本（未计自有资本的利息）且未扣无工资报酬的自有劳动。随后是两条会使损害被\textbf{高估}的 caveats：年度波动容不下全部 adaptation（``farmers will undertake these adaptations only if the benefits exceed the costs''\citep[第~381 页]{deschenes2007economic}，所以偏误方向明确），以及 CO$_2$ 浓度上升的增产效应未入模型（第~381 页）。

后续文献指出的：\citet{merel2021climate} 在第~1207 页把这条批评线概括为 panel approach 依赖 weather fluctuations 而非 climatic differences，因而可能反映不了 long-run adaptation；D-G 在 \S I.B 的自认（第~359 页）正是其源头。\citet{hsiang2016climate} 把这套设计的前提形式化为 marginal treatment comparability assumption（第~6 页，见 \chref{ch:hsiang-2016}）；\chref{ch:merel-gammans-2021} 推导出漏掉 penalty term 的面板系数相对 long-run response 的偏差，并给出同时估计 short-run 与 long-run response 的增广设定。

本讲义的评注：D-G 的地位不主要在它的点估计，而在三件事——把 identification 从截面挪到组内、把利润口径的测度缺陷交代到底、以及把「面板测不到全部 long-run 调整」写成了论文自己的结论。读这篇论文要同时拿到这三件；漏掉第三件，就看不出方法线下一步（\chref{ch:burke-hsiang-miguel-2015}、\chref{ch:merel-gammans-2021}）所回应的问题从哪里来。

\section{在本领域脉络中的位置}

论文回答的旧问题：截面 Ricardian 路线的 omitted variables bias 严重到什么程度（答案：hedonic 估计对口径选择极敏感，不可靠）。它开的新问题：panel approach 识别的 short-run response 与 long-run effect 差多少。后来的文献从三个方向处理这个差距：Burke 与 Emerick（2016）的 \index{long differences}long differences（时间线 2016）、Hsiang（2016）对 marginal treatment comparability 的部分检验（\chref{ch:hsiang-2016}）、Mérel 与 Gammans（2021）的 climate penalty 增广回归（\chref{ch:merel-gammans-2021}）。本讲义的归纳：在 \chref{ch:info-climate-change-adaptation} 的方法线上它是 panel 线的第一站，其自认局限（第~359 页）是该线此后各站的出发点。

% 章末参考文献（只列本章引用的条目），勿删
\chapterrefs
